\runninghead{Themes and Movement}
\section{The Propers: Themes and Movement}
\zlabel{triptych:concise:themes:start}%

The Church opens the liturgical year without the \work{Gloria in excelsis} (GIRM~53), setting seven Missal prayers and antiphons shared across Years~A, B, and~C beside five Year~B readings and chants from the Order of Readings for Mass (\work{Ordo lectionum Missae}) on the First Sunday of Advent in Year~B. Three starkly different scriptural voices meet the assembly. Isaiah asks why God has let Israel wander and hardened its heart, begs him to tear open the heavens, and admits that no one rises up to take hold of him (Isa.~63:16b--17, 19b; 64:2b--7; Vg.~63:16b--17; 64:1, 3b--8).\footnote{English Scripture quotations use the public-domain Douay--Rheims (Challoner) study translation rather than the approved English proclaimed from the U.S.\ \work{Lectionary for Mass}. Two numbering shifts affect these passages: in Isaiah, Hebrew and U.S.\ Lectionary 63:19b and 64:2b--7 correspond to Vulgate and Douay--Rheims 64:1 and 64:3b--8 (shifting chapter~64 by one verse; Vulgate Isaiah~64 citations are marked ``Vg.'' throughout); in the Psalms, citations give Hebrew and U.S.\ Lectionary numbers first, with Vulgate numbers in brackets or parentheses (Ps.~25[24], 80[79], 85[84]). Historical orientation for every appointed Scripture passage appears on the preceding page, \emph{Scriptural Date and Location}.} St.~Paul then thanks God that the Corinthians lack no grace while awaiting Christ's unveiling (1~Cor.~1:3--9), before Jesus on the Mount of Olives leaves his servants under one parting command: ``Watch'' (Mark~13:33--37).

\subsection{Introductory Rites: Trust in Waiting}

At the entrance, the antiphon \emph{Ad te levavi} (Ps.~25[24]:1--3) prays, ``To thee, O Lord, have I lifted up my soul... for none of them that wait on thee shall be confounded.'' In Latin, \emph{qui sustinent te} (``those who wait on thee'') means holding fast under pressure while help tarries. St.~Augustine (\work{Enarrationes in Psalmos}, Ps.~24 §§1--3) hears Christ praying in the person of the Church to reverse Adam's prideful self-trust, which left man fragile enough to fear death ``from the smallest insect.'' St.~Robert Bellarmine (\work{Explanatio in Psalmos}, Ps.~24, explanations~1--3) describes the believer finding ``briers and thorns'' in creatures and clinging to God's promises against human and demonic mockery.

From the Penitential Act, the Mass moves directly to the Collect, \emph{Da, quaesumus, omnipotens Deus} (``Grant, we beseech thee, almighty God''; the modern Latin and English prayer is not printed here). For historical comparison alone, an ancient Roman Postcommunion prayer (\work{Postcommunio}~683) with the same opening words in the \work{Sacramentarium Gelasianum Vetus} (Book~II, lxxxiv; its exact lexical alignment with the modern 2002/2008 Collect was not collated offline) asks almighty God to grant his whole family the will (\emph{voluntatem}) to go forth fittingly in righteous deeds (\emph{in operibus iustis aptos occurrere}) to meet Christ at his coming and inherit the heavenly kingdom.

\subsection{Liturgy of the Word: Lament, Endowment, and Vigilance}

The First Reading omits three blocks from the communal lament of Isaiah~63:7--64:12~Vg.\ (Heb.~63:7--64:11): 63:16a (Abraham and Israel not acknowledging the speakers), 63:18--19a~Heb.\ (Vg.~63:18--19, enemies trampling the sanctuary), and 64:1--2a~Heb.\ (Vg.~64:2--3a, fire kindling brushwood and the dread clause ``When thou shalt do wonderful things, we shall not bear them''). As a result, the wish ``O that thou wouldst rend the heavens, and wouldst come down'' (64:1~Vg.; Heb.~63:19b) stands directly beside the past-tense ``thou didst come down, and at thy presence the mountains melted away'' (64:3b~Vg.; Heb.~64:2b). From unseen blessings (64:4~Vg.) and God meeting ``him that rejoiceth, and doth justice'' (64:5~Vg.), the lament descends to stained deeds (``as the rag of a menstruous woman,'' 64:6~Vg.) and God's hidden face (64:7~Vg.), ending as clay before the Potter (64:8~Vg.).

The Responsorial Psalm (Ps.~80[79]:2ac, 3b, 15--16, 18--19; R.~v.~4) omits the tribal names of Joseph, Ephraim, Benjamin, and Manasseh (vv.~2b, 3a), the middle lament of tears and the wild boar tearing the vine (vv.~5--14), and verse~17's imprecation. The appointed verses pass straight from calling on the Shepherd upon the Cherubim to asking him to visit his vineyard, confirm ``the son of man'' (vv.~15--18), and restore the people's voice: ``we will call upon thy name'' (v.~19), where Isaiah~64:7~Vg.\ (Heb.~64:6) lamented that none calls on him. Where a U.S.\ parish sings either permitted Advent Common Psalm (Ps.~25[24] or Ps.~85[84], Lectionary No.~174; GIRM~61), it prays psalms already used at the Entrance, Gospel Acclamation, and Communion.

In the Second Reading (1~Cor.~1:3--9), Paul's thanksgiving sits immediately before his rebuke of Corinthian factions (1:10--12) and spiritual immaturity (3:1--3). Yet he thanks God that the Corinthians are enriched in utterance and knowledge and lack no grace while awaiting the \emph{revelatio} (Greek \emph{apokalupsis}, ``unveiling'') of Jesus Christ (1:7), promising that the faithful God (\emph{Fidelis Deus}) who calls them into his Son's fellowship will confirm them ``unto the end without crime'' (1:8--9).

Outside Lent, the \emph{Alleluia} stays before the Gospel (GIRM~62). Its verse, \emph{Ostende nobis, Domine} (Ps.~85[84]:8, ``Shew us, O Lord, thy mercy; and grant us thy salvation''), repeats the psalm refrain's verb \emph{ostende} (``show''), asking for what Augustine (\work{Enarr.}, Ps.~84 §6) calls ``Thy Christ.'' Mark~13:33--37 opens right after omitted verse~32 (``neither the angels in heaven, nor the Son, but the Father''). Jesus compares himself to a householder traveling abroad who gives his servants authority over every work and commands the doorkeeper (\emph{janitor}) to watch (v.~34); turning directly to the disciples, he commands them to stay awake because they do not know when the master comes --- evening, midnight, cockcrow, or morning (v.~35) --- and ends with a charge to all: ``Watch'' (v.~37). The Profession of Faith follows (Nicene Creed, or Apostles' Creed permitted, GIRM~68).

\subsection{Liturgy of the Eucharist and Concluding Rites: Gifts and Fruit}

No proper Offertory antiphon stands in the Missal, leaving that chant to local selection (GIRM~74). The modern Prayer over the Offerings (\emph{Suscipe, quaesumus, Domine}, ``Accept, we beseech thee, O Lord'') is not printed here; an ancient Roman counterpart in the \work{Sacramentarium Veronense} (XXVIIII.372, closely matching the 2002 digital Missal copy, with printed-edition bounds noted in the scope appendix) offers gifts drawn ``from thy bestowed benefits'' (\emph{de tuis... collata beneficiis}) so that temporal devotion may become eternal redemption.

The Order of Mass (\work{Ordo Missae}) provides two Prefaces of Advent (GIRM~364; their modern wording is not printed here), paired with Eucharistic Prayer~I (which has no proper Advent insert), II, or~III, excluding Eucharistic Prayer~IV (GIRM~365a--d). Advent thanksgiving looks across Christ's two comings (St.~Cyril of Jerusalem, \work{Catechesis}~15; Augustine on Ps.~25[24]:10) and St.~Bernard of Clairvaux's middle coming (\work{Sermones in Adventu}~3.4, 5.1). At Communion, the antiphon \emph{Dominus dabit benignitatem, et terra nostra dabit fructum suum} (Ps.~85[84]:13, ``The Lord will give goodness: and our earth shall yield her fruit'') returns to Psalm~85(84) (GIRM~87). The Prayer after Communion (\emph{Prosint nobis}, ``May they profit us,'' not printed here) closes the Communion Rite, followed by the Concluding Blessing --- ordinary blessing, or an optional solemn blessing or prayer over the people --- and the dismissal (GIRM~90).

\subsection{Three Interpretive Paths}

Catholic commentators answered the hardest verses of these passages in distinct ways; following their compatible arguments leads to three whole-Mass interpretations: \emph{Divine Forbearance and the Unveiling of an Already-Present Lord}, asking why God seems silent while hearts harden and the master tarries; \emph{The Potter's Clay, the Vine, and the Faithfulness of God}, asking how stained, divided believers can hope to stand blameless without overriding human freedom; and \emph{Vigilant Stewardship of the Household Between the Two Advents}, asking whom Mark~13 commands to watch across every century and what concrete life counts as watching.
\zlabel{triptych:concise:themes:end}
